\newglossaryentry{mpc}{
  name={Многосторонние конфиденциальные вычисления},
  description={криптографический подход, позволяющий нескольким участникам совместно вычислять значение функции на приватных входных данных без раскрытия этих данных друг другу}
}

\newglossaryentry{protocol}{
  name={Протокол},
  description={формализованная последовательность действий и сообщений, выполняемых участниками для достижения заданного результата вычисления}
}

\newglossaryentry{secret-sharing}{
  name={Разделение секрета},
  description={метод представления секрета в виде набора долей, при котором отдельная доля не раскрывает исходное значение, а восстановление возможно только при наличии достаточного числа долей}
}

\newglossaryentry{garbled-circuit}{
  name={Зашифрованная схема},
  description={представление вычисляемой функции в виде булевой схемы, элементы которой заменены случайными метками и зашифрованными таблицами истинности}
}

\newglossaryentry{homomorphic-encryption}{
  name={Гомоморфное шифрование},
  description={вид шифрования, позволяющий выполнять операции над шифротекстами так, что после расшифрования получается результат соответствующих операций над открытыми данными}
}

\newglossaryentry{oblivious-transfer}{
  name={Ослеплённая передача},
  description={криптографический протокол, в котором получатель выбирает одно из сообщений отправителя, не раскрывая свой выбор, а отправитель не узнаёт, какое сообщение было получено}
}

\newglossaryentry{adversary-model}{
  name={Модель противника},
  description={набор предположений о возможностях злоумышленника, включая возможность пассивного наблюдения, отклонения от протокола или преждевременного прекращения вычисления}
}

\newglossaryentry{fairness}{
  name={Справедливость (fairness)},
  description={свойство протокола, при котором результат вычисления получают либо все честные участники, либо ни один из них}
}

\newglossaryentry{robustness}{
  name={Устойчивость (robustness)},
  description={способность протокола сохранять корректность и завершать вычисление при допустимом моделью числе отказавших или действующих злоумышленно участников}
}

\newglossaryentry{tee}{
  name={Доверенная среда выполнения (Trusted Execution Environment, TEE)},
  description={изолированная аппаратно-программная среда, защищающая исполняемый код и обрабатываемые данные от остальной части системы}
}

\newglossaryentry{oram}{
  name={Скрытый доступ к памяти (Oblivious Random Access Memory, ORAM)},
  description={метод организации обращений к памяти, скрывающий от наблюдателя адреса и последовательность операций чтения и записи}
}

\newglossaryentry{beaver-triples}{
  name={Тройки Бивера (Beaver triples)},
  description={предварительно сформированные секретно разделённые случайные значения $a$, $b$ и $c$, для которых выполняется равенство $c=ab$; используются для защищённого умножения секретных величин}
}

\newglossaryentry{mpc-alliance}{
  name={MPC Alliance},
  description={международное отраслевое объединение организаций, содействующее распространению, внедрению и стандартизации технологий многосторонних конфиденциальных вычислений}
}